\chapter{Higher-Order Android Malware Analysis with \hgann{}}
\label{chap:hgann}

% Drafted under T-021 from notes/chapter5-source-records/. Governing decisions:
% D-027 and D-058 the Drebin binary corpus and empirical matrix; D-028 no ablation, the confound
% is declared and Chapter 1 carries no isolation claim; D-029 the individual contribution;
% D-008 the split of hypergraph formalism between Chapter 2 and this chapter;
% D-010 the qualified duplication wording. Revised under T-005 and D-046 from
% notes/chapter5-source-records/mathematical-revision-proposal.md: restored construction
% algorithm, operator analysis, complexity bounds, metric-consistency findings.
% The five figures are the author's own, supplied on 11 September 2026 (T-006, D-047); see
% figures/hgann-mal/README.md. They replace the earlier TikZ pipeline and hyperedge diagrams.

\section{Research Problem and Relation to Hybroid}
\label{sec:hgann:1}

This chapter answers RQ3, stated in Section~\ref{sec:intro:rqs}: \rqThreeText{} It takes up
gap~\gapref{3}. What that gap records is an asymmetry in the evidence. The cost of projecting a
group relation onto pairwise edges has been measured \cite{kang2024low, wang2024graphs}, and the
conditions under which the projection loses nothing are known \cite{agarwal2006higherorder,
chitra2019randomwalks}, yet the benefit of avoiding it in Android malware analysis rests on a
single published study \cite{zhang2023android}. The weighting usually applied alongside the
higher-order encoding carries its own unsettled evidence \cite{fang2025kaa}.

\hgann{} represents an application as a hypergraph whose vertices are the methods of its
function-call graph. Three families of hyperedge are laid over those vertices. The first records
proximity in the call graph, the second groups methods whose learned representations are close,
and the third follows a backward program slice from a call site that reaches a security-sensitive
API. The three are merged into one incidence structure, a hypergraph operator with learned
attention over hyperedges reduces the application to a single vector, and a classifier with two
heads decides malicious against benign and, separately, the class within a labelled taxonomy.
The source publication numbers the pipeline in eleven stages \cite{norouzian2025hgannmal}, and the
chapter keeps those numbers. Stages 0 to 2 catalogue and unpack the packages and recover the call
graph and the slices, and Stage 3 computes the node features. Stages 4 to 6 build and assemble the
hypergraph, Stage 7 forms mini-batches, Stage 8 is the attention operator, and Stages 9 and 10 are
the read-out and the two heads. The construction is described in Sections~\ref{sec:hgann:4}
to~\ref{sec:hgann:8} and the evaluation in Sections~\ref{sec:hgann:9} to~\ref{sec:hgann:13}.

Three claims are made for the work, and each is stated here in the form the evidence supports.
The first is the construction: a method-level function-call graph is lifted to a hypergraph by
three rules of different kinds, one topological, one semantic and one derived from a program
slice, with a weight on the third that rises with the number of data-transmission calls the slice
contains. The second is the operator: membership in a hyperedge is weighted by a learned score
rather than by a binary indicator, so the aggregation is trained rather than fixed by the
incidence structure. The third is the comparison, which measures the resulting system against two
pairwise graph networks and two hypergraph operators without attention, on two corpora and three
tasks, under one protocol.

Two limits belong here rather than in the limitations section, because they qualify what is being
claimed rather than what was found. The comparison changes the hyperedge vocabulary and the
aggregation together, so a difference in the totals cannot be assigned to either change on its
own; Section~\ref{sec:hgann:13} states what experiment would separate them and
Section~\ref{sec:hgann:14} carries the consequence. And the field this work enters is narrow. A
sweep of the published record finds one other hypergraph-based Android malware detector
\cite{zhang2023android}, so the comparison against hypergraph baselines is a comparison against
operators taken from other domains rather than against competing Android systems.

\subsection*{Relation to \hybroid{}}

Chapter~\ref{chap:hybroid} and this chapter answer one question about context in two ways, and
they are not stages of a single development. \hybroid{} widens the evidence base: it observes the
application twice, once as a program and once as a source of network traffic, and asks what the
combination adds. \hgann{} keeps to one source of evidence and widens the relation instead: it
asks what is gained when a group of methods can be recorded as a group rather than as a set of
pairs. The first buys context by paying for a second observation, with the acquisition cost that
Chapter~\ref{chap:hybroid} reports. The second buys context inside an observation already made.
Chapter~\ref{chap:discussion} takes up that trade directly, and it is the reason the two chapters
sit next to each other.

Their results are not comparable, and no comparison between them is offered anywhere in this
dissertation. The corpora differ, the class definitions differ, and the protocols differ. Where
the source publication of this chapter cites \hybroid{} as third-party prior work on hybrid
analysis, it is naming the candidate's own earlier work, and this chapter states the relationship
rather than repeating that framing.

The work was published at ISC 2025 \cite{norouzian2025hgannmal}, the candidate first of two
authors with his supervisor as the second. The method, its implementation, the experiments and the
manuscript are the candidate's own.
Section~\ref{sec:intro:publications} states the same in the form the affidavit covers.

\section{Related Work}
\label{sec:hgann:2}

Section~\ref{sec:background:graphs} defined the representations and
Section~\ref{sec:background:learning} the operators. What this section adds is the systems: which
of them have been built for Android, what they represent, and where the work of this chapter sits
among them.

\subsection*{Graph representations of Android applications}

Detection over program graphs predates graph neural networks. Adagio embedded call graphs with a
neighbourhood hash kernel and classified the result \cite{gascon2013structural}, and MaMaDroid
abstracted call sequences to package or family level before modelling transitions between them
\cite{onwuzurike2019mamadroid}. Their learned successors keep the graph and train the
representation over it. The survey that organises this line covers control-flow graphs, call
graphs, dependency graphs, heterogeneous application-and-API graphs, system-call graphs and
network-flow graphs, and no hypergraph operator appears among them \cite{bilot2024survey}.

One of those successors is close enough to this chapter to be worth stating in detail. MsDroid
gives up the whole program and works on local snippets around security-sensitive APIs, one
subgraph per snippet, classifying the subgraphs and combining their verdicts; it locates malicious
snippets under application-level labels alone, and it returns an explanation naming the control
flows responsible \cite{he2023msdroid}. The localisation is the same instinct that produces the
slice-aware hyperedges of Section~\ref{sec:hgann:7}, reached through a different construct.
MsDroid cuts subgraphs out of the call graph, while this chapter keeps one graph and lays a
hyperedge over the methods a slice touches.

What every member of the line shares is the restriction Section~\ref{sec:background:graphs}
describes. An edge relates two methods. A behaviour realised by several methods acting together
is representable only as a set of pairs, and the grouping is not recoverable from that set.

\subsection*{Hypergraph learning}

Learning over hypergraphs supplies the missing construct. HGNN defined a spectral convolution over
the incidence structure \cite{feng2019hypergraph}, its successor generalised the propagation with
a random-walk transition matrix and is the operator this chapter calls HGNN+ and the source
publication calls HGNNP \cite{gao2023hgnn}, and attention was carried across by making membership
itself real-valued and learned \cite{bai2021hypergraphattention}. The line has since moved again:
a stability-oriented successor from the same group supersedes the operator used as a baseline here
\cite{gao2026hgnnv2}, which is where a comparison run today would start.

\subsection*{Hypergraphs for Android malware, and how narrow that field is}

A sweep of the published record finds exactly one prior application of hypergraph learning to
Android malware, and it is worth describing precisely, both because it is the only true competitor
and because the source publication of this chapter misattributes it once in its own text
\cite{zhang2023android}. That work builds a function-call graph at method granularity, describes
each vertex by six centrality measures together with a one-hot encoding of what the function does,
and forms hyperedges by two rules: the $K$-order call neighbours of a function, and the set of
functions sharing a permission. It evaluates HGNN, HGNN+, HyperGCN and UniGNN against GCN, GIN and
GraphSAGE, on Malnet-Tiny and on Drebin, and it reports its best Drebin result as 97.1 per cent
accuracy with HGNN+. Its hyperedge weights are fixed in advance and no learned weighting is
evaluated anywhere in it, so every member of a hyperedge contributes equally to the aggregation.

\hgann{} differs from that work in three respects. It adds a third construction rule that neither
topology nor permissions supply, taking a hyperedge from the backward slice of a
security-sensitive call site, so that a group of methods jointly implicated in reaching a
sensitive API becomes one object. It weights that family by the transmission and access calls the
slice contains, which is the only place in the construction where a hyperedge carries a weight
other than one. And it replaces fixed membership with a learned score, which is the change RQ3
asks about. Being the second entrant in a two-member field is itself worth stating plainly: the
hypergraph baselines in Section~\ref{sec:hgann:10} are general operators applied to this
chapter's own representation, because there is almost nothing else to compare against.

\subsection*{The comparison that the published record does not support}

A reader will want the two Android hypergraph systems placed side by side, and the published
figures do not allow it. Zhang et al.\ report 97.1 per cent accuracy for HGNN+ on the
twenty-family Drebin task over 4,664 samples under an 80/20 split \cite{zhang2023android}, while
Section~\ref{sec:hgann:11} reports 95.0 per cent for \hgann{} and 93.3 per cent for the same
operator on the same task over the same samples under a 70/20/10 split. Setting those numbers
against each other would compare implementations, feature sets, splits and training regimes at
once.

The clearest evidence that such a comparison misleads is contained in the pair itself. One
operator, HGNN+, on one corpus and one task, differs by 3.8 points of accuracy between the two
papers, and that difference is larger than three of the four margins this chapter reports over its
own baselines. The problem is not peculiar to these two publications. A recent study that
re-implemented twelve Android detectors under one toolchain names unfair comparison as a
field-wide defect, on the ground that published approaches are evaluated on datasets of different
sizes, different ratios of benign to malicious samples and different training and test proportions
\cite{liu2026unraveling}, which is the same objection in general form \cite{arp2022dos,
daoudi2021lessons}. This chapter therefore reports what its own protocol measures and makes no
claim against the published figures of the other work.

\subsection*{What the evaluation literature requires}

Three further lines bear on the chapter and are taken up where they apply rather than here.
Methodological work on machine learning in security names the pitfalls that this chapter's
evaluation has to be checked against \cite{arp2022dos}, and studies of detector behaviour under
obfuscation, drift and adversarial perturbation give the conditions under which published accuracy
figures fail to hold \cite{gao2024comprehensive, zhang2020apigraph, liu2026unraveling}. Attacks
aimed specifically at graph-based Android detection show what an adversary can do to the
representation itself \cite{zhao2021structural, song2026fcghunter, li2025efficient}.
Section~\ref{sec:hgann:14} states what each of them implies for the results reported here.

\section{Observation and Threat Model}
\label{sec:hgann:3}

Everything this chapter measures follows from one decision about what is observed. \hgann{} reads
the application package and does not run it. The manifest and the Dalvik bytecode are the whole
of the evidence, which fixes both the cost of the analysis and the limits on what it can see.

The cost is low and it scales. No device, no emulator and no scripted interaction is needed, and
the analysis of one application is independent of every other, so the pipeline parallelises across
a corpus without coordination. The second modality of Chapter~\ref{chap:hybroid} depends on
traffic that the CICAndMal2017 authors captured on physical handsets, and
Section~\ref{sec:discussion:3} compares the two observation requirements directly.

The limits are the standard limits of static analysis on Android, stated in
Section~\ref{sec:background:android:analysis}, and three of them bear on the representation this
chapter builds. Code loaded at runtime through the class loader is absent from the recovered
bytecode, so no vertex exists for it. Reflective invocation resolves a target from a string at
runtime, so the call edge is missing from a statically recovered call graph unless the analyser
models the reflection explicitly, which the pipeline described here does not. Packing and native
payloads move the interesting behaviour outside the DEX files that are parsed at all. Each of the
three removes vertices or edges from the object the model consumes, and a representation that is
built on relations between methods is affected by a missing method more than a representation
built on a bag of permissions.

The adversary assumed by the evaluation is not adaptive. The corpora are historical collections,
their samples were written without knowledge of this detector, and no sample is modified in
response to what the detector learns. That assumption is standard for the reported task and it is
also the assumption Chapter~\ref{chap:advml} argues should be stated rather than left implicit,
which is why it is stated here. It has a consequence the chapter does not evade: an adversary who
rewrites the application changes the call graph the extraction recovers, and everything downstream
of that graph inherits the change. Section~\ref{sec:hgann:14} works that exposure out against the
published attacks on function-call-graph detectors.

\section{APK Processing and Function-Call Graph Extraction}
\label{sec:hgann:4}

Stage 0 builds a catalogue. Every application package under a root directory is hashed
with SHA-256, and the hash is joined against the label list distributed with the corpus, so each
sample carries a benign label or a family label before any analysis runs. Working from the hash
rather than from the file name matters for a corpus assembled from several sources, since it makes
duplicate files visible; Section~\ref{sec:hgann:9} returns to what was and was not done with that
information.

In Stage 1 each package is decompressed with \texttt{apktool} and the Python \texttt{zipfile}
module, which yields the manifest and every \texttt{classes*.dex} file the package contains.
Applications built with multi-dex are therefore covered, and the bytecode of all of them enters
the next stage. The reported runs used apktool, MobSF, Androguard with its DAD decompiler and a
ready-made code2vec model.

Androguard \cite{desnos2012android}, which the source runs beside MobSF in Stage 2a, disassembles
the DEX files and produces the object the rest of the chapter is built on: a function-call graph
$G = (V, E)$ at method granularity, in which a
vertex is a method identified by its class, its name and its descriptor, and a directed edge
records that one method may invoke another. The same pass emits a flat list of pairs that records,
for each method, which framework APIs it calls. That list is what allows a flagged API to be
located at its call sites rather than merely counted, and Section~\ref{sec:hgann:5} uses it for
exactly that.

Two properties of this object were fixed in Section~\ref{sec:background:graphs:pairwise} and are
not restated here beyond their consequences. A statically recovered call graph
over-approximates, so an edge in $G$ records that an invocation is possible and not that it
occurs. And the recovered graph depends on the analyser: a different toolkit, or a different
treatment of the component lifecycle and of callbacks registered with the framework, yields a
different vertex and edge set for the same application. Both properties travel through every
construction in Section~\ref{sec:hgann:7}, because each of the three hyperedge rules is defined
over $G$ or over an artefact derived from it.

\section{Security-Sensitive API Identification and Program Slicing}
\label{sec:hgann:5}

Stage 2 decides which parts of an application are worth isolating. In Stage 2a two tools run over
each package. MobSF parses the manifest, disassembles the bytecode and reports the permissions an
application requests together with the framework APIs it uses. Of the permissions reported, only
those whose status is \emph{dangerous} in the platform's own classification are retained, which
narrows attention to the capabilities that guard user data and device functions
\cite{aosp2026permissions}.

The APIs reached from those capabilities are then sorted into two classes by a hand-written rule.
An \emph{access} API reads something the user would consider private, and the examples the source
publication names are contacts, location, short messages and storage. A \emph{transmission} API
moves data off the device, with HTTP clients, raw sockets and \texttt{sendTextMessage} as its
examples. The pairing is deliberate. Data exfiltration is a sequence in which something is read
and then sent, and a rule that marks both ends of that sequence gives the construction in
Section~\ref{sec:hgann:7} something to join into one group. The source publication describes the
classification as a heuristic.

Stage 2b then computes a backward slice for every call site of a marked API. The variables that reach
the call site are identified first, and the control-flow graph of the enclosing code is walked
backwards, keeping only the blocks and statements that influence those variables. What the stage
stores is a compact control-flow graph per slice. The purpose is localisation rather than
verification: the slice does not establish that data flows from a source to a sink, in the sense
that a taint analysis would \cite{arzt2014flowdroid}, and it is not offered as evidence of
exfiltration. It marks a neighbourhood of code that participates in reaching a sensitive call, and
that neighbourhood becomes one hyperedge.

\section{Node Features}
\label{sec:hgann:6}

In Stage 3 each vertex of the call graph receives a feature vector assembled from four groups,
listed in Table~\ref{tab:hgann:features} and drawn in Figure~\ref{fig:hgann:node-features}. The
groups are concatenated and the result is standardised to zero mean and unit variance across the
corpus, giving the feature matrix $X$ with one row per method. The figure labels the blocks
$h_{\text{op}}$, $s_{\text{struct}}$, $a_{\text{api}}$ and $c_{\text{2v}}$, and the last of them is
the embedding written $c_v$ in Section~\ref{sec:hgann:7}. With $O$ the opcode vocabulary, the
width is $F = \lvert O \rvert + 328$.

\begin{sidewaysfigure}
  \centering
  \includegraphics[width=\textheight]{hgann-mal/hgann-mal-node-features.pdf}
  \caption[Construction of the node feature vector in Stage 3]{Construction of the node feature
  vector in Stage~3. For each method $v$ of the function-call graph, an opcode histogram, four
  structural centrality metrics, four suspicious-API flags and counts, and a 320-dimensional
  code2vec embedding are computed from the artefacts of Stages~1 to~2b. The four blocks are
  concatenated into $x_v$ and z-score normalised, which yields the node-feature matrix
  $X \in \mathbb{R}^{\lvert V \rvert \times F}$ consumed by the hypergraph constructor in
  Stage~6.}
  \label{fig:hgann:node-features}
\end{sidewaysfigure}

\begin{table}
  \caption{The four groups of node feature. The width of the first group is the size of the
  opcode vocabulary.}
  \label{tab:hgann:features}
  \centering
  \begin{tabular}{llp{58mm}}
    \toprule
    Group & Width & What it records \\
    \midrule
    Opcode histogram   & vocabulary   & the distribution of Dalvik instructions in the method body \\
    Structural metrics & 4            & in-degree, out-degree, PageRank and betweenness centrality of the method in $G$ \\
    API indicators     & 4            & whether the method reaches an access API and whether it reaches a transmission API, with the counts of each within the slices attached to it \\
    Semantic embedding & 320          & a code2vec representation of the decompiled method body \\
    \bottomrule
  \end{tabular}
\end{table}

The first group describes what a method does at instruction level, and it is the feature the
program-code branch of Chapter~\ref{chap:hybroid} also rests on, there as a learned embedding of
the opcode sequence rather than as a histogram. Operand values are not used, which removes the
constants and register numbers that vary between builds of the same code. The histogram can be
taken as raw counts or weighted by inverse document frequency across the corpus.

The second group places the method in the call graph rather than describing its body, and it is
the only group that changes when the graph changes. The third connects the vertex to the analysis
of Section~\ref{sec:hgann:5}, so that a method reaching a sensitive API is distinguishable before
any hyperedge is built. Both are small, and together they contribute eight of the dimensions.

The fourth group is the largest. code2vec represents a method as a bag of paths through its
abstract syntax tree and learns a fixed-width vector from that bag \cite{alon2019code2vec}. The
published method consumes Java source, and the released model was trained on a corpus of Java
source, whereas the input here is Dalvik bytecode recovered from a distributed package. The
candidate confirms that the two were bridged by decompilation: each method was decompiled to Java
with DAD, the decompiler built into Androguard, and embedded with a ready-made, pre-trained
code2vec model. The source publication records neither step, and the width it reports differs
from the released model's.

The embedding of a method therefore depends on how completely the decompiler recovers its source,
as well as on the model. That dependence matters for the construction too. The semantic hyperedges
of Section~\ref{sec:hgann:7} are built from cosine similarities between these embeddings, so the
decompiler and the model shape one of the three construction rules as well.

\section{Hyperedge Construction}
\label{sec:hgann:7}

The lifting from a graph to a hypergraph is where the representation is decided. Stages 4A to 4C
apply three rules of different kinds to the same vertex set, and Stages 5 and 6 merge and assemble
what they produce. Write $\mathcal{H} = (\mathcal{V}, \mathcal{E}, w)$ for the result, with
$\mathcal{V} = V$ the methods of the call graph. Figure~\ref{fig:hgann:hyperedge-generation} draws
one hyperedge of each kind together with the parameters of its rule. Its panels keep the names
$H_{\mathrm{K}}$, $H_{\mathrm{SEM}}$ and $H_{\mathrm{API}}$ that the source gives the three lists,
which Equation~\ref{eq:hgann:merge} writes as the candidate sets $\mathcal{C}_{\text{topo}}$,
$\mathcal{C}_{\text{sem}}$ and $\mathcal{C}_{\text{sus}}$.

\begin{sidewaysfigure}
  \centering
  \includegraphics[width=\textheight]{hgann-mal/hgann-mal-hyperedge-generation.pdf}
  \caption[Hyperedge generation in Stage 4]{Hyperedge generation in Stage~4: three complementary
  constructions over one vertex set of methods. (4A)~Topological hyperedges take the $K$-hop ball
  around every method in the call graph, with $K = 3$ and uniform weight $w = 1$. (4B)~Semantic
  hyperedges are $k$-nearest-neighbour groups in the 320-dimensional code2vec space, with $k = 5$,
  admitted only above the similarity floor $\tau \approx 0.80$. (4C)~Slice-aware hyperedges are
  built from the backward slices of Stage~2b, extended with shortest paths of length at most $L_p$
  between slice members, and weighted $w = 1 + 0.2\,\#\text{transmit} + 0.1\,\#\text{access}$.
  A candidate with fewer than 2 or more than 64 members is rejected. Node colour marks hyperedge
  membership only.}
  \label{fig:hgann:hyperedge-generation}
\end{sidewaysfigure}

\paragraph{Topological hyperedges.} The source collects, for each method $v$, its callers and
callees up to order $K$. Let $\vec{d}_G(v,u)$ be the number of call edges on a shortest directed
path from $v$ to $u$, with $\vec{d}_G(v,u) = \infty$ when there is none, and let
$\operatorname{dist}_G(v,u) = \min\{\vec{d}_G(v,u), \vec{d}_G(u,v)\}$. Then
\begin{equation}
  e_{\text{topo}}(v) = \{\, u \in V : \operatorname{dist}_G(v, u) \le K \,\}, \qquad K = 3 ,
  \label{eq:hgann:topo}
\end{equation}
holds $v$ together with its callers and callees up to $K$ steps away. This convention follows the
wording of the source. A distance measured with call directions ignored would also admit the other
callees of a caller and would give larger hyperedges, and the source does not record which
convention the implementation used.
This is the rule the earlier hypergraph work on Android uses, there in combination with a
permission-based rule \cite{zhang2023android}. It records a neighbourhood as one object rather
than as a set of pairs, which is the whole of the higher-order claim at this point: a vertex and
its $K$-hop neighbourhood enter the operator together.

\paragraph{Semantic hyperedges.} A nearest-neighbour search over the code2vec embeddings, computed
with the FAISS library, groups methods whose representations are close. Write
$c_v \in \mathbb{R}^{320}$ for the embedding of method $v$, the fourth group of
Table~\ref{tab:hgann:features}; the full feature vector $x_v$ also carries the other three groups,
and the search uses the embedding alone. The candidate confirms that the search compares cosine
similarities, $\operatorname{sim}(a, b) = a^{\top} b / (\lVert a \rVert\, \lVert b \rVert)$, and
that it runs within one application. With $\mathcal{N}_k(v)$ the $k$ methods of the same
application, other than $v$, that score highest under $\operatorname{sim}$,
\begin{equation}
  e_{\text{sem}}(v) = \{v\} \cup \{\, u \in \mathcal{N}_k(v) : \operatorname{sim}(c_v, c_u) \ge \tau \,\},
  \qquad k = 5, \quad \tau \approx 0.80 .
  \label{eq:hgann:sem}
\end{equation}
The original algorithm, whose test panel~4B of Figure~\ref{fig:hgann:hyperedge-generation}
reproduces, names the search output $D$, a letter that suggests a distance, and keeps a neighbour
when $D \ge \tau$. The test agrees with the cosine similarity, since for a distance the inequality
would run the other way. The query
method is included and is not counted among its $k$ neighbours, so
$\lvert e_{\text{sem}}(v) \rvert \le k + 1$; the source does not say whether its search counted the
query itself. Nothing in this rule refers to the call graph, which is the point of including it:
two methods that never call one another, and that sit in different parts of the application, are
joined when they do the same kind of work. The source names polymorphic malware and code reuse as
the intended cases. Because the search stays inside one application, a semantic hyperedge joins
methods of that application only, and code shared across a family can make the hyperedges of
different applications alike without joining them.

\paragraph{Slice-aware hyperedges.} Each backward slice $s$ from Section~\ref{sec:hgann:5} yields
one hyperedge over the methods the slice contains, extended by the methods on short call paths
between them. Let $F_s$ be that member set, fixed before any path is added. For a directed path $P$
in the call graph, let $V(P)$ be its methods and $\ell(P)$ its number of call edges, and let
$P_{uv}$ be a shortest path from $u$ to $v$. Then
\begin{equation}
  e_{\text{sus}}(s) = F_s \;\cup
  \bigcup_{\substack{u, v \in F_s,\ u \ne v\\ P_{uv} \text{ exists},\ \ell(P_{uv}) \le L_p}}
  V(P_{uv}) .
  \label{eq:hgann:sus}
\end{equation}
Only pairs of original members are joined, so a method that one path adds does not start further
paths. The original algorithm loops over the pairs of the set it is enlarging, and freezing $F_s$ is
the convention adopted here to make the result independent of the loop order. A pair with no
directed path contributes nothing. Where several shortest paths exist, the one chosen decides which
methods join. This family carries the only non-unit weight in the hypergraph:
\begin{equation}
  w(e_{\text{sus}}(s)) = 1 + 0.2\,\#\text{transmit}(s) + 0.1\,\#\text{access}(s) ,
  \label{eq:hgann:weight}
\end{equation}
so a slice that both reads private data and sends it weighs more than one that only reads, and the
transmission end weighs twice the access end. The weight is a construction heuristic, so it cannot
be read as a probability of exfiltration.

\paragraph{Merging.} Each rule proposes candidate triples $(e, w, t)$ of a member set, a weight and
a type tag. Collected as a multiset and passed through the post-processing of the source, they give
\begin{equation}
  \mathcal{C} = \mathcal{C}_{\text{topo}} \uplus \mathcal{C}_{\text{sem}} \uplus \mathcal{C}_{\text{sus}} ,
  \qquad
  \mathcal{E} = \operatorname{Post}_{\eta}\bigl( \{\, (e, w, t) \in \mathcal{C} :
  2 \le \lvert e \rvert \le c_{\max} \,\} \bigr) ,
  \qquad c_{\max} = 64 .
  \label{eq:hgann:merge}
\end{equation}
The lower bound removes singletons, which carry no relation, and with them every slice that has a
single member, since such a slice has no pair to join. The upper bound is described in the source
publication as necessary for stable message passing. Its prose says that hyperedges are capped,
whereas its algorithm rejects a candidate outside the interval. This chapter follows the algorithm:
truncation would need a rule for choosing which members to keep, and none is given. Rejection has a
consequence of its own. A method whose $K$-step neighbourhood holds more than 64 methods loses its
topological hyperedge entirely, so the topological family omits exactly the methods with the
largest neighbourhoods.

The operator $\operatorname{Post}_{\eta}$ removes duplicates. Two candidates with the same members
can carry different weights or tags, as a slice hyperedge and a topological one over the same
methods would, and the policy $\eta$ that decides which weight and which tag survive is not
recorded; this chapter does not supply one. Filtering can also leave a method in no hyperedge, for
instance an isolated method that belongs to no slice and has no semantic neighbour above $\tau$.
Its degree is then zero and the normalisation of Section~\ref{sec:hgann:8} is undefined for it, and
the source does not say whether such methods were dropped or handled in another way.

Each surviving hyperedge is tagged with its family in one byte, and the incidence matrix $H$ and
the diagonal weight matrix $W$ are assembled from the result and serialised per application.
Algorithm~\ref{alg:hgann:construction} restores the construction procedure of the source
publication with the conventions above written into it. It also adds the post-processing and the
sparse assembly, which the output line of the original announces and its body omits. The type tag
records something the operator never reads. The layer described in
Section~\ref{sec:hgann:8} consumes $H$, $W$ and $X$, so a hyperedge's family influences the model
only through its weight, and since two of the three families carry weight one, the topological and
semantic hyperedges are indistinguishable to the operator once they are built. That is a design
gap rather than an error, and Section~\ref{sec:hgann:13} names it as the cleanest available
extension.

\begin{algorithm}[tbp]
  \caption[Construction and assembly of the \hgann{} hypergraph]{Construction and assembly of the
  \hgann{} hypergraph. Stages 4A to 4C restore the algorithm of the source publication
  \cite{norouzian2025hgannmal} with the conventions of Equations~\ref{eq:hgann:topo}
  to~\ref{eq:hgann:sus} written in; stages 5 and 6 add the post-processing and assembly that its
  prose describes.}
  \label{alg:hgann:construction}
  \small
  \textbf{Input:} call graph $G = (V, E)$; embeddings $c_v$; slices $s \in \mathcal{S}$ with
  member sets $F_s$; parameters $K$, $k$, $\tau$, $L_p$, $c_{\max}$; duplicate policy $\eta$\par
  \textbf{Output:} sparse incidence matrix $H$, weight matrix $W$, type tags $t$
  \begin{algorithmic}[1]
    \Statex \textbf{Stage 4A: topological hyperedges}
    \State $\mathcal{C} \gets \varnothing$ \Comment{multiset of candidates}
    \For{$v \in V$}
      \State $e \gets \{u \in V : \operatorname{dist}_G(v,u) \le K\}$
      \If{$2 \le \lvert e \rvert \le c_{\max}$}
        \State $\mathcal{C} \gets \mathcal{C} \uplus \{(e, 1, \texttt{topo})\}$
      \EndIf
    \EndFor
    \Statex \textbf{Stage 4B: semantic ($k$-NN) hyperedges}
    \State $(I, S) \gets \textsc{Faiss}(\{c_v\}_{v \in V}, k)$
      \Comment{$k$ others per method; $S$ holds cosine similarities}
    \For{$v \in V$}
      \State $e \gets \{v\} \cup \{I_{v,j} : S_{v,j} \ge \tau\}$
      \If{$2 \le \lvert e \rvert \le c_{\max}$}
        \State $\mathcal{C} \gets \mathcal{C} \uplus \{(e, 1, \texttt{sem})\}$
      \EndIf
    \EndFor
    \Statex \textbf{Stage 4C: slice-aware hyperedges}
    \For{$s \in \mathcal{S}$}
      \State $e \gets F_s$; \ $w \gets 1 + 0.2\,\#\text{transmit}(s) + 0.1\,\#\text{access}(s)$
      \For{$(u, v) \in F_s \times F_s$ with $u \ne v$} \Comment{$F_s$ is not enlarged}
        \If{$P_{uv}$ exists and $\ell(P_{uv}) \le L_p$}
          \State $e \gets e \cup V(P_{uv})$
        \EndIf
      \EndFor
      \If{$2 \le \lvert e \rvert \le c_{\max}$}
        \State $\mathcal{C} \gets \mathcal{C} \uplus \{(e, w, \texttt{sus})\}$
      \EndIf
    \EndFor
    \Statex \textbf{Stage 5: post-processing}
    \State $\mathcal{E} \gets \operatorname{Post}_{\eta}(\mathcal{C})$
      \Comment{policy $\eta$ not recorded}
    \Statex \textbf{Stage 6: sparse assembly}
    \For{$j = 1, \ldots, \lvert \mathcal{E} \rvert$, with $(e_j, w_j, t_j)$ the $j$-th element
      of $\mathcal{E}$}
      \State $H_{vj} \gets 1$ for every $v \in e_j$; \ $W_{jj} \gets w_j$; \ $t(j) \gets t_j$
    \EndFor
    \State \Return $H$, $W$, $t$
  \end{algorithmic}
\end{algorithm}

\section{Hypergraph Attention Model}
\label{sec:hgann:8}

The operator is defined over the structure fixed in
Section~\ref{sec:background:graphs:hypergraphs}, whose notation is used here without restatement:
the incidence matrix $H$, the weighted vertex degree $d(v)$ and the hyperedge degree $\delta(e)$
with their diagonal matrices $D_v$ and $D_e$, and the diagonal weight matrix $W$. Write
$n = \lvert \mathcal{V} \rvert$, and assume for now that $d(v) > 0$ for every method;
Section~\ref{sec:hgann:7} records that filtering can violate this. The hypergraph Laplacian is
\begin{equation}
  L = I - P_0 , \qquad P_0 = D_v^{-1/2} H W D_e^{-1} H^{\top} D_v^{-1/2} ,
  \label{eq:hgann:laplacian}
\end{equation}
and for a signal $f \in \mathbb{R}^{n}$ over the methods it has the energy
\begin{equation}
  f^{\top} L f = \frac{1}{2} \sum_{e \in \mathcal{E}} \frac{w(e)}{\delta(e)}
  \sum_{u, v \in e} \left( \frac{f_u}{\sqrt{d(u)}} - \frac{f_v}{\sqrt{d(v)}} \right)^{2} ,
  \label{eq:hgann:energy}
\end{equation}
where the inner sum runs over ordered pairs \cite{zhou2006learning}. The energy is small when the
degree-normalised signal varies little inside the constructed groups, and each hyperedge counts in
proportion to its weight and in inverse proportion to its size. It describes the operator and is
not a term of the training loss. One layer of hypergraph convolution propagates through $P_0$:
\begin{equation}
  X^{(l+1)} = \sigma\!\left( D_v^{-1/2} H W D_e^{-1} H^{\top} D_v^{-1/2} X^{(l)} \Theta^{(l)} \right) ,
  \label{eq:hgann:conv}
\end{equation}
with $\Theta^{(l)}$ the learned projection of layer $l$ and $\sigma$ a non-linearity. Read from
right to left, the operator gathers vertex states into the hyperedges that contain them, averages
within each hyperedge through $D_e^{-1}$, scales by the hyperedge weight, and returns the result to
the member vertices with the symmetric normalisation of $D_v^{-1/2}$. One layer therefore couples
every pair of methods that share a hyperedge. The matrix $P_0$ is positive semidefinite and similar
to the row-stochastic matrix $D_v^{-1} H W D_e^{-1} H^{\top}$, so its eigenvalues lie in $[0,1]$.
Equation~\ref{eq:hgann:conv} thus multiplies each eigencomponent of $L$ with eigenvalue $\lambda$ by
$1 - \lambda$, and the components with the highest energy are damped most. This is the smoothing
behind the low-pass bias recorded in Section~\ref{sec:background:learning:hypergraphconvolution}.

\subsection*{Attention over hyperedges}

Equation~\ref{eq:hgann:conv} weights a vertex's hyperedges by $W$, which is fixed when the
hypergraph is built. Attention replaces the binary membership indicator by a learned real value,
so that the strength with which a vertex participates in a hyperedge is trained
\cite{bai2021hypergraphattention}. The layer index is dropped from here on. Let
$X \in \mathbb{R}^{n \times F}$ hold the input states as rows, let
$\Theta \in \mathbb{R}^{F \times F'}$ be the projection, and set $Y = X\Theta$, with
$y_v \in \mathbb{R}^{F'}$ the row of method $v$ read as a column vector. A hyperedge is first
represented by the mean of the projected states of its members,
\begin{equation}
  z_e = \frac{1}{\delta(e)} \sum_{u \in e} y_u ,
  \label{eq:hgann:hedgerep}
\end{equation}
and the coefficient joining a vertex to one of its hyperedges is the normalised score
\begin{equation}
  \alpha(v, e) = \frac{\exp\bigl( \sigma( a^{\top} [\, y_v \,\Vert\, z_e \,] ) \bigr)}
                      {\sum_{e' \ni v} \exp\bigl( \sigma( a^{\top} [\, y_v \,\Vert\, z_{e'} \,] ) \bigr)} ,
  \label{eq:hgann:attn}
\end{equation}
with $a \in \mathbb{R}^{2F'}$ a learned vector and $\Vert$ concatenation. Collecting the
coefficients into $H_{\text{att}}$, whose entry at $(v,e)$ is $\alpha(v,e)$ where $H$ has a one and
zero elsewhere, gives the layer actually used:
\begin{equation}
  X' = \sigma\!\left( D_v^{-1/2} H_{\text{att}} W D_e^{-1} H_{\text{att}}^{\top} D_v^{-1/2}
              X \Theta \right) .
  \label{eq:hgann:attnprop}
\end{equation}
The degree matrices are those of the binary $H$, as in the source publication; degrees recomputed
from $H_{\text{att}}$ would define a different layer. Gradients reach $H_{\text{att}}$ as well as
$X$ and $\Theta$, so the membership weights are learned. The support is not. $H_{ve} = 0$ implies
$\alpha(v,e) = 0$, and a finite score gives every present membership a positive coefficient, so
training reweights the memberships the construction produced and cannot add a method to a
hyperedge or remove one.

Equations~\ref{eq:hgann:hedgerep} and~\ref{eq:hgann:attn} state the additive form of the score,
which the source publication leaves as an uninstantiated similarity function and which the
candidate confirms is what was implemented. The published formulation also writes the score
between two vertices while indexing it by a hyperedge; the form above resolves that by scoring a
vertex against a representation of the hyperedge, which is what the implementation computes.

\subsection*{The messages between methods}

Expanding Equation~\ref{eq:hgann:attnprop} shows what one method receives from another. The layer
runs in two stages. Each hyperedge first forms a message, and each method then combines the
messages of its hyperedges:
\begin{equation}
  \mu_e = \frac{1}{\delta(e)} \sum_{u \in e} \frac{\alpha(u,e)}{\sqrt{d(u)}}\, y_u ,
  \qquad
  x'_v = \sigma\biggl( \frac{1}{\sqrt{d(v)}} \sum_{e \ni v} w(e)\, \alpha(v,e)\, \mu_e \biggr) .
  \label{eq:hgann:messages}
\end{equation}
Substituting the first expression into the second gives one matrix of couplings between methods,
\begin{equation}
  X' = \sigma\bigl( P_{\text{att}} X \Theta \bigr) ,
  \qquad
  P_{\text{att}}(v,u) = \frac{1}{\sqrt{d(v)\, d(u)}} \sum_{e \ni v,u} \frac{w(e)}{\delta(e)}\,
  \alpha(v,e)\, \alpha(u,e) .
  \label{eq:hgann:coupling}
\end{equation}
Two methods interact only through the hyperedges they share, and each shared hyperedge adds one
term. In that term the coefficient of the sender and the coefficient of the receiver appear as a
product. The construction weight $w(e)$ is a separate factor, so the slice heuristic of
Equation~\ref{eq:hgann:weight} and the learned coefficients act on a message independently. The
coefficients depend on the layer input through Equations~\ref{eq:hgann:hedgerep}
and~\ref{eq:hgann:attn}, which makes $P_{\text{att}}$ a function of $X$; for a fixed input it is a
symmetric matrix.

The key $z_e$ and the message $\mu_e$ are different aggregations over the same members. The key is
an unweighted mean, whereas the message weights each member $u$ by $\alpha(u,e) / \sqrt{d(u)}$, so
attention changes the average inside a hyperedge as well as the weighting across the hyperedges of
a method. It does not change the sign of any term. Every member enters $\mu_e$ with a positive
coefficient, so before the non-linearity each output is the projection of a non-negative
combination of input states, and one layer cannot form a difference between two members of a
hyperedge. The high-frequency component that standard hypergraph convolution discards has been
characterised \cite{li2026highpass}. Attention changes the proportions of that average and keeps
every coefficient positive, so the propagation of one layer remains an averaging of member states.

Grouping the sum in Equation~\ref{eq:hgann:coupling} by hyperedge gives a structural statement:
\begin{equation}
  P_{\text{att}} = \sum_{e \in \mathcal{E}} \frac{w(e)}{\delta(e)}\, \psi_e \psi_e^{\top} ,
  \qquad
  (\psi_e)_v = \frac{\alpha(v,e)}{\sqrt{d(v)}} ,
  \label{eq:hgann:rankone}
\end{equation}
with $\alpha(v,e) = 0$ for $v \notin e$. Each hyperedge adds a positive semidefinite matrix of rank
one, and because every weight is positive, $P_{\text{att}}$ is positive semidefinite for every
input. The statement concerns the propagation matrix at a fixed input. It carries no guarantee about
the convergence of training.

Because $\alpha(v,e)$ depends on both the method and its hyperedge, two different hyperedge
groupings can produce the same fixed matrix $P_0$ yet different attentive propagation matrices.
The reported experiments do not determine whether the learned model used this capacity or whether
it caused the measured accuracy differences.

\subsection*{What the mechanism can and cannot be claimed to do}

Two further properties bound what the chapter claims for the operator.

The first concerns the ranking that the score induces. Splitting $a = (a_1, a_2)$ across the
concatenation writes the argument of $\sigma$ in Equation~\ref{eq:hgann:attn} as $q_v + b_e$, with
a query term and a hyperedge key
\begin{equation}
  q_v = a_1^{\top} y_v , \qquad b_e = a_2^{\top} z_e .
  \label{eq:hgann:split}
\end{equation}
If $\sigma$ is strictly increasing, then for a fixed layer input and any two hyperedges $e_1$ and
$e_2$ that contain $v$,
\begin{equation}
  \alpha(v, e_1) > \alpha(v, e_2) \iff b_{e_1} > b_{e_2} .
  \label{eq:hgann:ranking}
\end{equation}
The proof is short. Because $\sigma$ and $\exp$ are strictly increasing, the numerators of
Equation~\ref{eq:hgann:attn} are ordered as the keys are, and the two coefficients share one
positive denominator. Any two methods therefore rank the hyperedges they share in the same order,
and the keys alone fix that order. Methods can still attend to different hyperedges, since each
normalises over its own. A non-decreasing activation such as ReLU can merge two hyperedges into a
tie but cannot reverse their order between methods. The argument is the one Brody et al.\ give for
the additive score of pairwise graph attention \cite{brody2022gatv2}, applied to the score of
Equation~\ref{eq:hgann:attn}. It holds for each attention head separately and says nothing about
hypergraph attention with other scoring functions. Under either kind of activation, the restriction
on wording that the chapter keeps is thus a result for this operator: within one layer, \hgann{}
can learn which hyperedges matter, and it cannot reverse their order from one method to another.

The query term does not cancel, however. With $\sigma$ a leaky ReLU of slope 0.2 and keys $-1$
and $1$, the coefficients are approximately $0.23$ and $0.77$ for $q_v = 0$, and $0.12$ and $0.88$
for $q_v = 2$. The order agrees and the values differ. Cancellation is exact when $\sigma$ is affine
on the arguments of one method, as a leaky ReLU is when they share a sign.

The second concerns normalisation. Uniform coefficients do not reproduce
Equation~\ref{eq:hgann:conv}, because the degrees in Equation~\ref{eq:hgann:attnprop} are those of
the binary $H$. Let $r_v$ be the number of hyperedges containing $v$ and
$R = \operatorname{diag}(r_v)$. Equal scores give $H_{\text{att}} = R^{-1} H$ and hence
\begin{equation}
  P_{\text{att}} = R^{-1} P_0 R^{-1} ,
  \label{eq:hgann:uniform}
\end{equation}
which reduces to $P_0 / r^2$ when every method lies in $r$ hyperedges, and to $P_0 / 4$ in the
example above. An ablation that replaced the learned coefficients by uniform ones would therefore
change the scale of propagation as well as its weighting. A comparison with the non-attentive
operator needs a stated normalisation convention before a difference can be read as an effect of
attention.

\subsection*{Read-out and classification}

Applications differ in size, so the vertex states are collapsed into one vector per application.
Two variants are named in the source publication, an unweighted mean over the vertices and a
learned attention pooling.

Training and inference run on batches. The source assembles a batch of $B$ applications by placing
their incidence and weight matrices on a block diagonal and stacking their feature matrices,
\begin{equation}
  H_{\mathcal{B}} = \operatorname{diag}(H_1, \ldots, H_B) , \qquad
  W_{\mathcal{B}} = \operatorname{diag}(W_1, \ldots, W_B) , \qquad
  X_{\mathcal{B}} = \begin{bmatrix} X_1 \\ \vdots \\ X_B \end{bmatrix} ,
  \label{eq:hgann:batch}
\end{equation}
together with an index that maps each row of $X_{\mathcal{B}}$ to its application. No hyperedge
spans two blocks, so degrees, keys, coefficients and messages are all computed inside one
application. The output for an application therefore does not depend on the other applications in
its batch, and the read-out uses the index to pool each block on its own.

The layer is also equivariant under relabelling. Let $\Pi$ reorder the methods of an application
and $Q$ its hyperedges, and write $\Phi(X, H, W)$ for the output of
Equation~\ref{eq:hgann:attnprop}. Then
\begin{equation}
  \Phi(\Pi X, \Pi H Q^{\top}, Q W Q^{\top}) = \Pi\, \Phi(X, H, W) , \qquad
  g(\Pi X, \Pi H Q^{\top}, Q W Q^{\top}) = g(X, H, W) ,
  \label{eq:hgann:invariance}
\end{equation}
where $g$ is the application vector. The first identity holds because the degree matrices, the keys
and the coefficients all permute with $H$, and the second because both read-outs are symmetric in
the vertex states. The prediction therefore does not depend on the order in which method rows or
hyperedge columns are stored. This concerns the hypergraph once it is built. It gives no resistance
to obfuscation, and order-dependent choices earlier in the pipeline, such as tie-breaking among
nearest neighbours or among shortest paths, can change $H$ itself.

That vector passes through a shared multi-layer perceptron with two hidden layers, rectified
linear activations and dropout, and then into two heads. The first emits a single logit whose
sigmoid is the probability that the application is malicious. The second emits one logit per class
of the labelled taxonomy, normalised by a softmax. Training minimises the sum of the two
objectives,
\begin{equation}
  \mathcal{L} = \lambda_{\text{bin}} \operatorname{BCE}(z_{\text{bin}}, y_{\text{bin}})
              + \lambda_{\text{fam}} \operatorname{CE}(z_{\text{fam}}, y_{\text{fam}}) ,
  \label{eq:hgann:loss}
\end{equation}
weighted by two constants.

Equation~\ref{eq:hgann:loss} applies both terms unconditionally. On the corpus of
Section~\ref{sec:hgann:12}, benign
applications form one of the labelled classes and the second term is defined for them. On the
corpus of Section~\ref{sec:hgann:11} the classes are twenty malware families. A benign application
has no family among them, and neither has a malicious application from one of the smaller
families, so the family term has no target for either. One well-defined specification averages each
term over the training samples that carry its label,
\begin{equation}
  \mathcal{L} = \frac{\lambda_{\text{bin}}}{\lvert \mathcal{D}_{\text{bin}} \rvert}
  \sum_{i \in \mathcal{D}_{\text{bin}}} \operatorname{BCE}(z_{\text{bin},i}, y_{\text{bin},i})
  + \frac{\lambda_{\text{fam}}}{\lvert \mathcal{D}_{\text{fam}} \rvert}
  \sum_{i \in \mathcal{D}_{\text{fam}}} \operatorname{CE}(z_{\text{fam},i}, y_{\text{fam},i}) ,
  \label{eq:hgann:maskedloss}
\end{equation}
with $\mathcal{D}_{\text{bin}}$ and $\mathcal{D}_{\text{fam}}$ the non-empty sets of samples that
have a binary label and a class label.

\section{Datasets, Tasks, and Split Policy}
\label{sec:hgann:9}

Two corpora are used, and they answer to different needs. Drebin is the reference collection of
the field, which makes results on it comparable with a long line of published work, and its family
labels support a classification task that a category taxonomy cannot. CICMalDroid is newer, it
carries benign applications alongside malicious ones, and its labels are categories of behaviour
rather than families. Using both guards against a conclusion that holds only for one collection,
which is the reason the source publication gives for the second corpus.

\subsection*{Drebin}

Drebin is the malicious part of a larger collection gathered between August 2010 and October 2012
from the Google Play Store, from Chinese and Russian alternative markets, from Android websites and
malware forums, and from the whole of the Malgenome project \cite{arp2014drebin}. Every sample in
that collection was submitted to VirusTotal and judged by ten antivirus scanners; an application
was counted as malicious when at least two of the ten flagged it, applications labelled adware were
then removed, and 5,560 malicious samples remained. The family names are taken from the labels of
one of the ten scanners. The source publication of this chapter counts 179 families in the set
\cite{norouzian2025hgannmal}, a figure the dataset paper does not state. Family sizes are heavily
skewed, so the classification task uses the twenty largest, which cover 4,664 of the samples and
are drawn with their sizes in Figure~\ref{fig:hgann:class-distribution}(a). The 159 smaller
families hold the other 896 samples, 16.1 per cent of the set, and they are left out of the
family task only, since binary detection uses all 5,560 malicious samples.

\begin{figure}[tbp]
  \centering
  \includegraphics[width=\linewidth]{hgann-mal/hgann-mal-class-distribution.pdf}
  \caption[Class distribution of the two corpora]{Class distribution of the two corpora.
  (a)~The twenty largest Drebin families, the classes of the \familyclassification{} task. The
  159 smaller families, with 896 samples, take no part in that task and remain in binary
  detection. (b)~The five CICMalDroid 2020 categories of the published distribution
  \cite{mahdavifar2020dynamic}; Benign is the negative class of binary detection on this corpus.
  The long tail of Drebin is the reason the source publication reports an F1 score, which it
  labels Macro-F1, alongside accuracy.}
  \label{fig:hgann:class-distribution}
\end{figure}

Three properties of that construction bear on the results. The malicious label is a vote of two
scanners out of ten, which is a lower bar than a majority and admits samples the other eight
consider harmless. The ground truth for \familyclassification{} is one vendor's naming, which is a
convention rather than a taxonomy of behaviour. And the 123,453 benign packages the original
DREBIN work evaluated against are not part of what is distributed, which is why a negative class
had to be found elsewhere.

Binary detection needs a negative class, and Drebin supplies none, so the benign side was drawn at
random from AndroZoo \cite{allix2016androzoo}. AndroZoo collects Android applications by crawling
Google Play, the Chinese markets Anzhi and AppChina, and several smaller repositories; the crawl
began in late 2011 and has continued since, and the collection held more than three million
applications when it was first described. Every application in it has been submitted to
VirusTotal, and what the collection stores per application is the number of antivirus products
that flagged it, together with the compilation date of its bytecode, its size and the market it
came from.
About 9,470 benign applications were used against the full set of 5,560 malicious ones, giving a
corpus of roughly 15,030 in which malware is 37.0 per cent.

Turning that stored count into a benign label requires a threshold, and the threshold is a
decision rather than a property of the data. The AndroZoo authors report the flagged share at one
antivirus product and again at ten, and state that antivirus results have to be treated with care;
the two thresholds give very different populations. Section~\ref{sec:hgann:14} states what
follows from assembling the two classes from different collections.

\subsection*{CICMalDroid}

CICMalDroid 2020 was collected between December 2017 and December 2018 and holds five categories,
four malicious and one benign. Its benign class is not a separate acquisition: applications that
fell into none of the four malicious categories were taken as benign and then checked against
VirusTotal. Its published distribution, drawn in Figure~\ref{fig:hgann:class-distribution}(b),
holds 11,598 applications, which is not the number collected: 17,341 applications were
gathered, 13,077 executed successfully in the dataset authors' sandbox, and the labelled
distribution retains 11,598 after the analysis records of the remainder failed to parse
\cite{mahdavifar2020dynamic}. This chapter uses the published distribution, so it inherits a
selection made by dynamic analysis, and Section~\ref{sec:hgann:14} records what that implies for a
static detector. The categories are far more even than the Drebin families, since the largest
holds 3.1 times as many samples as the smallest.

\subsection*{Three tasks, and they are not interchangeable}

Binary detection separates malicious from benign, and it is reported on both corpora. On Drebin
the negative class is the AndroZoo sample described above; on CICMalDroid it is the 1,795 benign
applications of Figure~\ref{fig:hgann:class-distribution}(b) against the remaining 9,803.

The \familyclassification{} task assigns a malicious application to one of the twenty Drebin
families. No benign application takes part, and the task presupposes that the sample is malicious.

The \categoryclassification{} task assigns an application to one of the five CICMalDroid
categories, one of which is Benign. It therefore includes the detection decision within a wider
one.

The distinction is kept throughout this chapter and throughout the dissertation, because the three
tasks have different priors, different error costs and different trivial baselines, and reporting
them under one name would make the numbers look comparable when they are not.

The two corpora do not even agree on what counts as malicious. Drebin removed every application
its scanners labelled adware, while CICMalDroid keeps Adware as one of its four malicious
categories. A sample of the kind Drebin discarded is therefore a positive case on one corpus and
absent from the other, which is one more reason the figures in
Sections~\ref{sec:hgann:11} and~\ref{sec:hgann:12} are not read against each other.

\subsection*{Split policy}

Each corpus is divided at random into 70 per cent for training, 20 per cent for testing and 10 per
cent for validation, and the same division is used for every method in the comparison, so the
baselines and \hgann{} see the same samples. At that test fraction the smallest Drebin family, with
41 samples, places about eight samples in the test partition in expectation, so a per-family score
for the smallest classes would rest on a handful of decisions. Two properties of that policy are
stated here and their consequences are worked out in Section~\ref{sec:hgann:14}: the division is
random rather than ordered by time, so training and test samples come from the same period, and
no deduplication or hash-aware grouping was applied, although Drebin is known to contain
repackaged near-duplicates.

\section{Binary Detection Results}
\label{sec:hgann:10}

Binary detection asks one question of an application: malicious or not. It is reported on both
corpora, and the two experiments differ in a way that matters for reading the numbers. On
CICMalDroid the benign applications belong to the corpus, and the task uses the 1,795 benign
samples against the 9,803 malicious ones. On Drebin they do not, because Drebin holds malware
only, so the benign side was assembled separately as Section~\ref{sec:hgann:9} describes.
Table~\ref{tab:hgann:binary} reports both. Figures~\ref{fig:hgann:f1} and~\ref{fig:hgann:accuracy}
plot its F1 and accuracy columns together with those of the multiclass table in
Section~\ref{sec:hgann:11}, for all four task and corpus combinations.

\begin{table}
  \caption{Binary detection, in per cent, as published. The source labels the first column
  Macro-F1 and states no averaging convention. Section~\ref{sec:hgann:14} shows that the \hgann{}
  values on Drebin are positive-class values for the malicious class, and that several rows cannot
  follow one convention. Best value in each column in bold.}
  \label{tab:hgann:binary}
  \centering
  \begin{tabular}{lrrrr@{\hskip 2em}rrrr}
    \toprule
    & \multicolumn{4}{c}{Drebin and AndroZoo} & \multicolumn{4}{c}{CICMalDroid} \\
    \cmidrule(lr){2-5} \cmidrule(lr){6-9}
    Method & F1 & Prec. & Rec. & Acc. & F1 & Prec. & Rec. & Acc. \\
    \midrule
    GCN            & 92.2 & 92.2 & 91.5 & 92.8 & 90.6 & 90.6 & 90.6 & 90.9 \\
    GraphSAGE      & 92.1 & 92.1 & 92.5 & 93.3 & 91.0 & 91.1 & 91.1 & 91.5 \\
    \midrule
    HGNN           & 96.2 & 95.8 & 96.3 & 96.8 & 94.2 & 94.8 & 94.3 & 94.6 \\
    HGNN+          & 96.1 & 95.9 & 96.1 & 97.1 & 95.9 & 96.1 & 96.5 & \textbf{97.7} \\
    \hgann{}       & \textbf{97.8} & \textbf{97.8} & \textbf{97.6} & \textbf{98.3}
                   & \textbf{96.2} & \textbf{96.2} & \textbf{97.6} & 97.6 \\
    \midrule
    Majority class & & & & 63.0 & & & & 84.5 \\
    \bottomrule
  \end{tabular}
\end{table}

\begin{figure}[tbp]
  \centering
  \includegraphics[width=\linewidth]{hgann-mal/hgann-mal-f1.pdf}
  \caption[F1 of the five methods on both corpora and both tasks]{F1 on both corpora and both
  tasks, as published in Tables~\ref{tab:hgann:binary} and~\ref{tab:hgann:multi}. The source labels
  these values Macro-F1, and Section~\ref{sec:hgann:14} shows that several rows cannot follow one
  averaging convention. The vertical axis is truncated at 80 so that differences among the
  hypergraph methods remain visible. HGNNP is the source publication's name for HGNN+.}
  \label{fig:hgann:f1}
\end{figure}

\begin{sidewaysfigure}
  \centering
  \includegraphics[width=\textheight]{hgann-mal/hgann-mal-accuracy.pdf}
  \caption[Accuracy against the majority baseline and the strongest baselines]{Accuracy reported
  in Tables~\ref{tab:hgann:binary} and~\ref{tab:hgann:multi}, divided at the majority baseline, at
  GraphSAGE and at HGNN+, which the figure labels HGNNP after the source publication. Variability
  across runs and the isolated contribution of attention are not established by these reported
  results.}
  \label{fig:hgann:accuracy}
\end{sidewaysfigure}

Three readings follow, and they are of unequal strength.

The clearest is the gap between the pairwise and the higher-order representations. On Drebin every
hypergraph operator sits above every pairwise one, by between 3.5 and 5.5 points of accuracy and
between 3.9 and 5.7 points of F1, and the ordering holds on CICMalDroid as well.
Figure~\ref{fig:hgann:accuracy} shows the same step on all four combinations: HGNN+ exceeds the
best pairwise operator, GraphSAGE in every case, by 3.8 to 6.7 points of accuracy, and in each
combination this step is larger than the difference between \hgann{} and HGNN+. All four
baselines were re-implemented on this pipeline and consumed the same call graph and the same node
features, so the difference belongs to
the representation rather than to the input. That is the comparison RQ3 asks for first, and it is
the one the evidence supports best.

The second is the gap between \hgann{} and the two hypergraph operators without attention. On
Drebin it is 1.2 points of accuracy against the stronger of them, and on CICMalDroid it is not
there at all: HGNN+ holds the best accuracy at 97.7 against 97.6, while \hgann{} holds the better
recall at 97.6 against 96.5 and the better F1 at 96.2 against 95.9. A system that leads on three
columns of four, and trails by one tenth of a point on the fourth, is level with its baseline on
that corpus rather than ahead of it. The source publication reads the same pattern as a
precision-recall trade-off, which is one explanation among several; single runs without dispersion
cannot separate a trade-off from noise, and Section~\ref{sec:hgann:14} returns to that.

The third concerns the base rate, and it is the reason the last row of the table is there. A
constant predictor scores 63.0 per cent on the Drebin corpus, where benign applications are the
majority, and 84.5 per cent on CICMalDroid, where they are 15.5 per cent of the samples. Reported
accuracy stands 35.3 points above the trivial baseline in the first case and 13.1 points above it
in the second. Neither figure is diminished by this, but the second is much less impressive than
97.6 per cent sounds, and the source publication reports no such baseline.

\subsection*{Drebin confusion matrix and class balance}

The Drebin binary experiment evaluated \hgann{} on 1,112 malicious and 1,894 benign applications.
The classifier identified 1,085 malware samples and missed 27. It raised 24 false alarms among
the benign applications and classified 1,870 as benign. These 51 errors across 3,006 decisions
give a false-positive rate of 1.3 per cent.

The experiment used a corpus with 37.0 per cent malware, whereas deployments may contain a lower
share. A separate sensitivity calculation holds recall and the false-positive rate fixed while
reducing the malware share to 10 per cent. Under this assumption, precision falls from 97.8 to
about 89.5 per cent because benign applications outnumber malicious applications nine to
one. This calculation illustrates the effect of class balance under fixed rates, as discussed in
Section~\ref{sec:background:evaluation}.

\section{Malware-Family Classification on Drebin}
\label{sec:hgann:11}

The second task assigns a malicious application to the family it belongs to. It runs on the 4,664
Drebin samples of the twenty largest families, shown in
Figure~\ref{fig:hgann:class-distribution}(a), and no
benign application takes part. Results are in the left panel of Table~\ref{tab:hgann:multi}.

\begin{table}
  \caption{Multiclass classification, in per cent, as published. The Drebin task assigns a
  malicious application to one of twenty families; the CICMalDroid task assigns an application to
  one of five categories, of which Benign is one. The two tasks are not comparable with each other.
  Section~\ref{sec:hgann:14} identifies two rows whose three columns cannot follow one averaging
  convention.}
  \label{tab:hgann:multi}
  \centering
  \begin{tabular}{lrrrr@{\hskip 2em}rrrr}
    \toprule
    & \multicolumn{4}{c}{Drebin, 20 families} & \multicolumn{4}{c}{CICMalDroid, 5 categories} \\
    \cmidrule(lr){2-5} \cmidrule(lr){6-9}
    Method & F1 & Prec. & Rec. & Acc. & F1 & Prec. & Rec. & Acc. \\
    \midrule
    GCN            & 86.1 & 85.9 & 86.2 & 87.8 & 82.4 & 82.0 & 82.4 & 82.7 \\
    GraphSAGE      & 86.3 & 86.7 & 86.2 & 88.0 & 82.3 & 82.2 & 82.6 & 83.4 \\
    \midrule
    HGNN           & 92.0 & 92.1 & 92.4 & 92.0 & 88.7 & 88.8 & 88.8 & 89.0 \\
    HGNN+          & 92.3 & 92.5 & 92.5 & 93.3 & 89.3 & 89.5 & 89.8 & 90.1 \\
    \hgann{}       & \textbf{94.2} & \textbf{94.1} & \textbf{93.8} & \textbf{95.0}
                   & \textbf{94.0} & \textbf{94.1} & \textbf{94.0} & \textbf{94.2} \\
    \midrule
    Majority class & & & & 19.8 & & & & 33.7 \\
    \bottomrule
  \end{tabular}
\end{table}

The ordering repeats: pairwise operators lowest, non-attentive hypergraph operators four to five
and a half points of accuracy above them, and \hgann{} 1.7 points above the better of those. Against
the trivial baseline the margin is wide here, since the largest family covers 19.8 per cent of the
samples.

One property of the task deserves more caution than the totals invite. The twenty families range
from 925 samples down to 41, a ratio above twenty to one. Read as a macro average, which is how
the source labels it, an F1 of 94.2 per cent would require the smallest families to be classified
nearly as well as the largest, since every class contributes equally to that average.
Section~\ref{sec:hgann:14} shows that the printed precision and recall of the same row exclude any
reading in which the three columns share one convention. The classes are also one vendor's
labels rather than an agreed taxonomy, as Section~\ref{sec:hgann:9} records, so a disagreement
between the model and the label is not always an error of the model. The families are further
related to one another: several of the twenty are known to contain repackaged applications built from shared
code. Shared code gives the applications of one family similar semantic hyperedges, each built
inside its own application under Section~\ref{sec:hgann:7}, and it is also the condition under
which duplication inflates a measured score. Section~\ref{sec:hgann:14} returns to that threat.

\section{Malware-Category Classification on CICMalDroid}
\label{sec:hgann:12}

The third task assigns an application to one of the five CICMalDroid categories. Benign is one of
them, so this task subsumes the detection question rather than sitting beside it, and an
application misclassified between two malicious categories is a different kind of error from one
misclassified as benign. Results are in the right panel of Table~\ref{tab:hgann:multi}.

The margin over the non-attentive operators is 4.7 points of F1 and 4.1 of accuracy, the widest of
the four task and corpus combinations reported anywhere in this chapter. The narrowest, on the
same corpus, is binary detection, where Section~\ref{sec:hgann:10} found no margin at all. The
ranking of the four combinations does not depend on the measure. Over the strongest baseline, the
F1 margins in Figure~\ref{fig:hgann:f1} are 4.7 points here, 1.9 and 1.6 on the Drebin family and
binary tasks and 0.3 on CICMalDroid binary detection, which is also the order of the accuracy
margins. That pattern is worth stating plainly because it constrains what may be concluded:
whatever the model
adds, it adds most where the decision is between several classes and least where it is between
two. A representation that groups methods by shared behaviour would plausibly help most in exactly
that setting, since categories such as adware and riskware are separated by what an application
does with what it collects rather than by whether it collects anything. The observation is offered
as a reading of the pattern and not as a finding, since one run per configuration cannot support
more.

\section{Ablation, Efficiency, and Interpretability Analysis}
\label{sec:hgann:13}

\subsection*{The ablation that was not run}

\hgann{} differs from the hypergraph baselines in two ways at once. It builds a different set of
hyperedges, fusing a topological, a semantic and a slice-aware family where the baselines take
$k$-hop and permission-based ones \cite{zhang2023android}, and it weights membership by a learned
score instead of a fixed indicator. Both changes are present in every number in
Tables~\ref{tab:hgann:binary} and~\ref{tab:hgann:multi}, and no experiment in this chapter varies
one while holding the other. The measured differences are therefore differences between two
systems, and they do not establish that attention is what produces them.

Four experiments would settle it, and they are recorded here so that the gap is specified rather
than merely admitted. First, the baseline operator without attention, run over the identical fused
hyperedge set, which isolates the aggregation. Second, \hgann{} run over topological hyperedges
alone, which isolates the vocabulary. Third, a leave-one-out over the three hyperedge families,
which apportions the vocabulary's contribution among them. Fourth, a sweep over $K$, $k$, $\tau$,
$c_{\max}$, $L_p$ and the two weight constants, none of which is varied anywhere in the source
publication. The first two are the ones that bear on RQ3; the second two would tell a reader how
sensitive the construction is to choices made once and never revisited.
Equation~\ref{eq:hgann:uniform} adds a condition to the first: uniform coefficients also rescale
propagation, so that comparison has to fix a normalisation convention before its difference can be
attributed to attention.

The question is not idle. On pairwise graphs, over node-property benchmarks and against
non-attentive baselines, no unmodified attentive model has been shown to outperform all of those
baselines on any dataset or task tested \cite{fang2025kaa}, and the scoring functions of both the
original and the revised formulations realise only finitely many neighbour rankings at bounded
width \cite{brody2022gatv2, fang2025kaa}. Neither result was stated for hyperedges, as
Section~\ref{sec:background:learning:attention} records. The static-ranking argument carries over
to the score used here, as Equation~\ref{eq:hgann:ranking} shows, while the finite-ranking bound
has not been examined for hyperedge scores. What the two results establish is that the attribution
cannot be assumed, and the evidence needed to make it was not gathered.

\subsection*{Efficiency}

The stated procedure admits bounds. Let $n$ be the number of methods, $m$ the number
of retained hyperedges, $\lvert \mathcal{S} \rvert$ the number of slices and
$\nu = \operatorname{nnz}(H) = \sum_{e \in \mathcal{E}} \lvert e \rvert$ the number of stored
memberships. Algorithm~\ref{alg:hgann:construction} proposes at most one topological and one
semantic candidate per method and one candidate per slice, and post-processing only removes
candidates. Semantic hyperedges hold at most $k + 1$ methods under the convention of
Equation~\ref{eq:hgann:sem}, and the others at most $c_{\max}$. Hence
\begin{equation}
  m \le 2n + \lvert \mathcal{S} \rvert , \qquad
  \nu \le (c_{\max} + k + 1)\, n + c_{\max} \lvert \mathcal{S} \rvert
  = 70\, n + 64\, \lvert \mathcal{S} \rvert ,
  \label{eq:hgann:sizes}
\end{equation}
and the representation grows linearly in the numbers of methods and slices. For one attention head
with input width $F$ and output width $F'$, a sparse evaluation of Equation~\ref{eq:hgann:attnprop}
computes the projection, the keys, the scores and the two incidence products in
\begin{equation}
  T_{\text{layer}} = O(nFF' + \nu F') , \qquad
  M_{\text{layer}} = O\bigl(nF + (n + m)F' + \nu + FF'\bigr)
  \label{eq:hgann:cost}
\end{equation}
time and memory, where $m \le \nu / 2$ because every retained hyperedge has at least two members.
The bound follows the two stages of Equation~\ref{eq:hgann:messages} and never forms
$P_{\text{att}}$. That matrix can be far denser than $H$, since two methods are coupled whenever
they share a hyperedge: it can hold up to $\sum_e \lvert e \rvert^2 \le c_{\max} \nu$ non-zero
entries, and the cap bounds this density. Further heads and layers add their own terms.

Construction costs more. With $\lvert E \rvert$ call edges, breadth-first searches from every
method build the topological hyperedges in $O(n(n + \lvert E \rvert))$. An exhaustive search within
one application evaluates the similarity of $O(n^2)$ pairs of 320-dimensional embeddings, and since
the FAISS index is unspecified, no sub-quadratic bound can be claimed for it. The slice paths need
one breadth-first search from each member of each slice,
$O\bigl(\sum_s \lvert F_s \rvert (n + \lvert E \rvert)\bigr)$ in total. The cap is applied after a
candidate is built, so it does not limit the cost of candidates that are then rejected, and
centrality computation and the semantic embedding add costs outside these bounds. None of these
quantities is a measurement. The claim in the source publication that the approach is inexpensive
because it is static is true of the analysis. For the representation, the bounds above are the only
evidence, and those for construction are quadratic in the number of methods.

\subsection*{Interpretability}

The slice-aware hyperedges are the part of the construction that could carry an explanation. Each
one is anchored at a call site of a security-sensitive API and covers the methods a backward slice
passes through, so an attention weight on such a hyperedge points at a concrete region of code and
at the API that made it interesting. Reporting the highest-weighted hyperedges for a sample, with
the call sites they contain, would turn a score into something an analyst can check.
Equation~\ref{eq:hgann:ranking} makes that report compact: within one layer and head every method
ranks its hyperedges by the same keys $b_e$, so one ranked list of keys per application summarises
what the layer attends to.

That was not done. The source publication states interpretability as future work, and no attention
weight is inspected, visualised or reported anywhere in it. A single retrained model on one corpus would be enough: the weights of
one application, ranked, with the API call at the anchor of each hyperedge, would turn this
section from a description of what the design permits into a result about what the model attends
to. It is named here because it is the cheapest substantial improvement available to the design,
and because the construction was built in a way that makes it reachable.

\subsection*{One design gap, and it is the cleanest extension available}

Section~\ref{sec:hgann:7} records that every hyperedge is tagged with the family that produced it,
and that the operator never reads the tag. A hyperedge reaches the layer as a column of $H$ and an
entry of $W$, so a topological hyperedge and a semantic one over the same methods are the same
object to the model, and only the slice-aware family is distinguished at all, through a weight it
happens to carry. The type is available, it is stored, and it is discarded.

Two ways of spending it suggest themselves. The weight matrix could carry a learned scalar per
family, so that the training decides how much a topological grouping is worth against a semantic
one, which costs three parameters. Or the attention score of
Equation~\ref{eq:hgann:attn} could take the type as an input. Entered additively, the type would
only shift the keys $b_e$, and Equation~\ref{eq:hgann:ranking} shows that every method would still
rank its hyperedges in one shared order; a weighting that depends on the family and the vertex
together needs a score that is not additive, such as the revised form of Brody et
al.\ \cite{brody2022gatv2}. Neither was tried here, and the second is the more
interesting because it is the version of the mechanism that the confound in
Section~\ref{sec:hgann:14} leaves unexamined.

\section{Limitations and Threats to Validity}
\label{sec:hgann:14}

The results above are measurements on two historical corpora under one protocol. This section
states the limits of what they show and closes with a brief account of threats to validity.

\paragraph{The gain cannot be attributed to attention.} \hgann{} changes the hyperedge vocabulary
and the aggregation together, and no experiment separates the two (Section~\ref{sec:hgann:13}).
The evidence supports a claim about the complete system: it performs better than four baselines on
three of the four task and corpus combinations. Whether learned weighting of hyperedges produces
that gain is an open question.

\paragraph{Every figure is a single run.} No configuration was repeated, and the tables therefore
carry no standard deviation, confidence interval or significance test. Small margins are affected
most. A tenth of a point separates \hgann{} from HGNN+ on CICMalDroid accuracy, and that difference
cannot be interpreted. Even the 1.2 points on Drebin would need repeated runs before they could be
called reliable.

\paragraph{The evaluation covers favourable conditions only.} Both splits are random, so training
and test samples come from the same collection period, and the figures do not describe performance
on applications released later \cite{pendlebury2019tesseract}. Neither corpus was obfuscated, and
no adversary was present. Under obfuscation, concept drift and adversarial examples, none of twelve
representative Android detectors kept the performance it reached in the ideal setting, and one
call-graph-based detector scored 0.912 of F1 on one year's test set and 0.463 on another's
\cite{gao2024comprehensive}. Sections~\ref{sec:hgann:10} to~\ref{sec:hgann:12} thus report the
conditions that a comparison between representations needs. A deployment would meet harder ones.

\paragraph{Some behaviour lies outside the representation.} Code fetched at runtime and targets
resolved by reflection are absent from the recovered call graph, and so are payloads held in packed
archives or native libraries (Section~\ref{sec:hgann:3}). A representation built from relations
among methods loses more from a missing method than one built from a bag of permissions.

\paragraph{The baselines are general operators.} GCN, GraphSAGE, HGNN and HGNN+ were
re-implemented on this pipeline's own representation. They test whether the higher-order encoding
and attention help within one design, but they do not place \hgann{} against the Android detectors
that the field uses as reference points. The hypergraph baseline has since been superseded by a
stability-oriented successor from the same group \cite{gao2026hgnnv2}, so the comparison reflects
that line of work as it stood in 2022.

\paragraph{The threat model is non-adaptive.} No adversary in this evaluation knows that the
detector exists, although published attacks act on the object this pipeline consumes. The first
structural attack on graph-based Android detection reports over 90 per cent success in feature
space and up to 100 per cent in problem space against function-call-graph detectors under
white-box access \cite{zhao2021structural}. FCGHunter rewrites an application with
semantics-preserving operators and reaches a white-box success rate of 87.9 per cent against forty
detectors built from graph-summary features \cite{song2026fcghunter}. A4M reports an average of
85.17 per cent against seven detectors built with MaMaDroid, APIGraph and a graph neural network,
under a perturbation budget of five per cent \cite{li2025efficient}. None of these attacks was run
against \hgann{}. Every hyperedge of Section~\ref{sec:hgann:7} is derived from the call graph or
from quantities computed over it, so methods and edges injected into that graph also change the
hyperedges. Attention operates downstream of an extraction that the adversary controls. \hgann{}
therefore inherits the exposure of its representation, and in line with Chapter~\ref{chap:advml}
this chapter makes no robustness claim.

\paragraph{Threats to validity.} The result tables head their first column Macro-F1 without stating
a convention. For the \hgann{} Drebin binary row a macro reading is impossible: over 1,112
malicious and 1,894 benign test applications, macro precision 97.8 and macro recall 97.6 at an
accuracy that rounds to 98.3 would require at least 1,897 true negatives. That row is read as
positive-class values. Under positive-class, macro and support-weighted averaging alike,
$F_1 \le (P + R)/2$, and four rows exceed that mean by more than the 0.1 point that rounding
allows: GCN and HGNN in the Drebin binary panel, \hgann{} in the \familyclassification{} panel, and
GCN in the \categoryclassification{} panel. No single convention fits these rows, so the
multiclass columns are not read as macro averages; accuracy is unaffected. On the data side, the
benign Drebin applications come from AndroZoo, which began crawling in late 2011, while Drebin
covers August 2010 to October 2012, so a classifier may separate the two collections by build tools
or API levels that have nothing to do with malice \cite{arp2022dos, allix2016androzoo}. The source
publication records that about 49 per cent of Drebin samples are repackaged applications, and none
was removed. Near-duplicates across the partitions may inflate the family scores of
Section~\ref{sec:hgann:11} more than the binary ones \cite{zhao2021impact, irolla2018duplication}.
CICMalDroid, for its part, keeps only applications that ran in its authors' sandbox
\cite{mahdavifar2020dynamic}, which may remove evasive samples that a static detector most needs
to catch.

\section{Chapter Summary}
\label{sec:hgann:15}

This chapter answered RQ3 by lifting a method-level function-call graph to a hypergraph, joining
its vertices by three rules of different kinds, and learning how strongly each method participates
in each group rather than fixing that strength when the structure was built. The system was
measured against two pairwise graph operators and two hypergraph operators without attention, on
Drebin and on CICMalDroid, for binary detection and for classification into families and into
categories.

Two findings hold. The higher-order representation is better than the pairwise one on every task
and corpus reported, by between three and eleven and a half points of accuracy depending on the
task and on which pair is compared, with the input to both held identical. The attention-weighted system is ahead of the non-attentive hypergraph
operators on three of four combinations and level on the fourth, and its widest margin appears on
the multiclass task rather than on detection. The analysis of Section~\ref{sec:hgann:8} sharpens
what the second finding can mean without adding evidence for it. Attention reweights the
memberships that the construction fixed, and within one layer it ranks the hyperedges of every
method in one shared order. A constructed example shows that it can still separate structures that
the fixed operator identifies, although nothing shows that training exploits this.

One thing does not hold, and the chapter states it rather than working around it. Because the
hyperedge vocabulary and the weighting were changed together, the improvement over the hypergraph
baselines cannot be attributed to attention, and the experiment that would separate them was not
run. Chapter~\ref{chap:discussion} takes the representational question up against
Chapter~\ref{chap:hybroid}, and Chapter~\ref{chap:conclusion} records the gap as the first item of
future work.
